% !PRINTED 55
\weberpage{55}

%<<CONT
Flächenstückes liegen, wobei Nullpunkte und Unstetigkeitspunkte $m$ter Ordnung wie $m$ solche Punkte erster Ordnung zu zählen sind.
Hieraus folgt, daß im Inneren eines endlichen Flächenstückes auch nur eine endliche Anzahl von Nullpunkten liegt.

In gleichem Sinne ergibt sich

7. Das Integral
\[
\frac{1}{2\pi i}\int u\,d\log\varphi(u)
\]
ist gleich der Summe der Werte von $u$, für die $\varphi(u)$ verschwindet, vermindert um die Summe der Werte von $u$, für die $\varphi(u)$ unendlich wird.

8. Eine Funktion, die im Endlichen gar keine Unstetigkeitspunkte besitzt, heißt eine \gesperrt{ganze Funktion}. \gesperrt{Eine ganze Funktion, die auch im Unendlichen endlich bleibt, ist eine Konstante.} Eine ganze Funktion, die auch im Unendlichen keine wesentlich singuläre Stelle hat, für die es also einen Exponenten $m$ derart gibt, daß $u^{-m}\varphi(u)$ für $u = \infty$ nicht unendlich wird, ist eine rationale Funktion.

9. Eine ganze Funktion, deren Reziproke gleichfalls ganz ist, heißt eine \gesperrt{Einheitsfunktion}. Eine solche Funktion wird im Endlichen weder Unendlich noch Null. Ihr Logarithmus ist daher ebenfalls eine ganze Funktion, und es folgt daraus, daß jede Einheitsfunktion in der Form $e^{g(u)}$ dargestellt werden kann, worin $g(u)$ eine ganze Funktion ist. Umgekehrt ist jeder Ausdruck von dieser Form eine Einheitsfunktion.

10. Nach den grundlegenden Untersuchungen von \gesperrt{Weierstrass} über die eindeutigen analytischen Funktionen (Abhandlungen der Berliner Akademie von 1876) gibt es immer eine ganze Funktion $G(u)$, die in den Unstetigkeitspunkten einer Funktion $\varphi(u)$ und nur in diesen verschwindet, und diese Funktion $G(u)$ ist durch die Unstetigkeitspunkte von $\varphi(u)$ bis auf eine Einheitsfunktion als Faktor bestimmt. Es ist dann $\varphi(u)G(u) = G_1(u)$ gleichfalls eine ganze Funktion, und man kann also jede Funktion $\varphi(u)$ als Quotienten zweier ganzen Funktionen
\[
\varphi(u) = \frac{G_1(u)}{G(u)}
\]
darstellen, worin $G(u)$ und $G_1(u)$ keine gemeinschaftlichen Nullpunkte haben.
